% GENERATED FILE. Edit canonical lessons, then run npm run generate:books.
% canonical-source-hash: fnv1a64:daad97f570459571
% canonical-lessons: LA-C113-nam, LA-C113-enim, LA-C113-ideo, LA-C113-itaque, LA-C113-causa

\chapter{For, For, For that reason, And so, Reason}
\label{ch:la-for-for-for-that-reason-and-so-reason}
\hlchaptermodality{113}

\begin{chapteropening}
\textbf{By the end of this chapter:} \emph{I can say for, for, for that reason, and so and a reason.}

Complete the last lesson of chapter 113: I can say for, for, for that reason, and so and a reason.
\end{chapteropening}

\section[nam]{nam — for (giving a reason)}

\label{lesson:LA-C113-nam}



\begin{quote}
Before the new one: say the Latin for why?, then the Latin for why?.
\end{quote}

\subsection*{You'll want to know: nam}
\textbf{nam} — \textquotedblleft{}for (giving a reason)\textquotedblright{}.

\begin{cousinweb}
\textbf{nam} is usually explained as an old case form of a pronoun stem \emph{na-}, \textquotedblleft{}that\textquotedblright{} (Lewis \& Short): \textquotedblleft{}for that\textquotedblright{}.
\end{cousinweb}

\begin{grammarlens}[title={What you've built}]
One more word for this chapter.
\end{grammarlens}

\subsection*{Your turn}
\begin{itemize}
\raggedright
  \item \emph{Say it:} \emph{nam}
  \item \emph{Say it:} \emph{nam}, once more
  \item \emph{Say it:} say \emph{nam}
  \item \emph{Recall:} say the Latin for why?, then the Latin for why?, then say \emph{nam} again
\end{itemize}

\begin{tcolorbox}[breakable,colback=teal!4,colframe=teal!35!black,title={Before you move on}]
What does nam mean? (\textquotedblleft{}For (giving a reason)\textquotedblright{}.) Say it once more.
\end{tcolorbox}
\section[enim]{enim — for, you see}

\label{lesson:LA-C113-enim}



\begin{quote}
Before the new one: say the Latin for why?, then the Latin for for.
\end{quote}

\subsection*{You'll want to know: enim}
\textbf{enim} — \textquotedblleft{}for, you see\textquotedblright{}.

\begin{cousinweb}
\textbf{enim} is often explained as a short pointing element followed by \emph{nam}, \textquotedblleft{}for\textquotedblright{} (Lewis \& Short), but its make-up is uncertain.
\end{cousinweb}

\begin{grammarlens}[title={What you've built}]
One more word for this chapter.
\end{grammarlens}

\subsection*{Your turn}
\begin{itemize}
\raggedright
  \item \emph{Say it:} \emph{enim}
  \item \emph{Say it:} \emph{enim}, once more
  \item \emph{Say it:} say \emph{enim}
  \item \emph{Recall:} say the Latin for why?, then the Latin for for, then say \emph{enim} again
\end{itemize}

\begin{tcolorbox}[breakable,colback=teal!4,colframe=teal!35!black,title={Before you move on}]
What does enim mean? (\textquotedblleft{}For, you see\textquotedblright{}.) Say it once more.
\end{tcolorbox}
\section[ideō]{ideō — for that reason}

\label{lesson:LA-C113-ideo}



\begin{quote}
Before the new one: say the Latin for for, then the Latin for for.
\end{quote}

\subsection*{You'll want to know: ideō}
\textbf{ideō} — \textquotedblleft{}for that reason\textquotedblright{}.

\begin{cousinweb}
\textbf{ideō} is built from \emph{id}, \textquotedblleft{}that\textquotedblright{}, and \emph{eō}, \textquotedblleft{}for that\textquotedblright{} (Lewis \& Short): for that reason.
\end{cousinweb}

\begin{grammarlens}[title={What you've built}]
One more word for this chapter.
\end{grammarlens}

\subsection*{Your turn}
\begin{itemize}
\raggedright
  \item \emph{Say it:} \emph{ideō}
  \item \emph{Say it:} \emph{ideō}, once more
  \item \emph{Say it:} say \emph{ideō}
  \item \emph{Recall:} say the Latin for for, then the Latin for for, then say \emph{ideō} again
\end{itemize}

\begin{tcolorbox}[breakable,colback=teal!4,colframe=teal!35!black,title={Before you move on}]
What does ideō mean? (\textquotedblleft{}For that reason\textquotedblright{}.) Say it once more.
\end{tcolorbox}
\section[itaque]{itaque — and so, therefore}

\label{lesson:LA-C113-itaque}



\begin{quote}
Before the new one: say the Latin for for, then the Latin for for that reason.
\end{quote}

\subsection*{You'll want to know: itaque}
\textbf{itaque} — \textquotedblleft{}and so, therefore\textquotedblright{}.

\begin{cousinweb}
\textbf{itaque} is \emph{ita}, \textquotedblleft{}so, thus\textquotedblright{}, with \emph{-que}, \textquotedblleft{}and\textquotedblright{} (EtymDB 2.1): and so.
\end{cousinweb}

\begin{grammarlens}[title={What you've built}]
One more word for this chapter.
\end{grammarlens}

\subsection*{Your turn}
\begin{itemize}
\raggedright
  \item \emph{Say it:} \emph{itaque}
  \item \emph{Say it:} \emph{itaque}, once more
  \item \emph{Say it:} say \emph{itaque}
  \item \emph{Recall:} say the Latin for for, then the Latin for for that reason, then say \emph{itaque} again
\end{itemize}

\begin{tcolorbox}[breakable,colback=teal!4,colframe=teal!35!black,title={Before you move on}]
What does itaque mean? (\textquotedblleft{}And so, therefore\textquotedblright{}.) Say it once more.
\end{tcolorbox}
\section[causa, causae]{causa, causae — a reason, a cause}

\label{lesson:LA-C113-causa}



\begin{quote}
Before the new one: say the Latin for for that reason, then the Latin for and so.
\end{quote}

\subsection*{You'll want to know: causa, causae}
\textbf{causa, causae} — \textquotedblleft{}a reason, a cause\textquotedblright{}.

\begin{cousinweb}
\textbf{causa} was spelled \emph{caussa} in older Latin (EtymDB 2.1); where it comes from is unknown. English \emph{cause} and \emph{because} come from it.
\end{cousinweb}

\begin{grammarlens}[title={What you've built}]
One more word for this chapter.
\end{grammarlens}

\subsection*{Your turn}
\begin{itemize}
\raggedright
  \item \emph{Say it:} \emph{causa}
  \item \emph{Say it:} \emph{causa}, once more
  \item \emph{Say it:} say \emph{causa}
  \item \emph{Recall:} say the Latin for for that reason, then the Latin for and so, then say \emph{causa} again
\end{itemize}

\begin{tcolorbox}[breakable,colback=teal!4,colframe=teal!35!black,title={Before you move on}]
What does causa, causae mean? (\textquotedblleft{}A reason, a cause\textquotedblright{}.) Say it once more.
\end{tcolorbox}
